\begin{trapbox}
	\begin{description}[style=nextline, leftmargin=2.8em, labelsep=0.5em, itemsep=0.35em]
		\item[\textbf{\textcolor{CTrap}{易错点一（线段混淆）}}] 误将割线定理记为 $PA \cdot AB = PC \cdot CD$ 或 $PA \cdot AB = PC \cdot PD$ 等形式。
		\item[\textbf{\textcolor{CTrap}{正确理解（乘积形式）：}}] 定理中等号两边都是“近段 $\times$ 远段”，即 $PA \cdot PB = PC \cdot PD$，其中 $A,C$ 是近交点，$B,D$ 是远交点。
		\item[\textbf{\textcolor{CTrap}{错因分析（对应不清）：}}] 对定理中线段对应关系不清晰，死记符号导致错误。
	\end{description}

	\medskip
	\begin{description}[style=nextline, leftmargin=2.8em, labelsep=0.5em, itemsep=0.35em]
		\item[\textbf{\textcolor{CTrap}{易错点二（内外混淆）}}] 当点 $P$ 在圆内时，仍然使用割线定理。
		\item[\textbf{\textcolor{CTrap}{正确理解（位置区别）：}}] 割线定理仅适用于点 $P$ 在圆外；当 $P$ 在圆内时，应使用相交弦定理（两弦相交于圆内），结论形式类似但条件不同。
		\item[\textbf{\textcolor{CTrap}{错因分析（条件忽略）：}}] 未区分定理的使用场景，将形式相同的结论套用错误条件。
	\end{description}

	\medskip
	\begin{description}[style=nextline, leftmargin=2.8em, labelsep=0.5em, itemsep=0.35em]
		\item[\textbf{\textcolor{CTrap}{易错点三（比例写反）}}] 在利用相似推导时，将对应边比例写反，如得到 $PA/PB = PC/PD$ 的错误关系。
		\item[\textbf{\textcolor{CTrap}{正确理解（正确比例）：}}] 正确比例应为 $PA/PD = PC/PB$，交叉相乘得 $PA \cdot PB = PC \cdot PD$。
		\item[\textbf{\textcolor{CTrap}{错因分析（对应顶点）：}}] 相似三角形的对应顶点没有找准，导致比例错误。
	\end{description}
\end{trapbox}
